\documentclass[11pt]{article}

\usepackage[utf8]{inputenc}
\usepackage[T1]{fontenc}
\usepackage{lmodern}
\usepackage{geometry}
\usepackage{amsmath,amssymb,amsfonts,amsthm,mathtools}
\usepackage{bm}
\usepackage{enumitem}
\usepackage{microtype}
\usepackage{hyperref}
\usepackage{titlesec}
\usepackage{booktabs}
\usepackage{array}
\usepackage{setspace}
\usepackage{mathrsfs}
\usepackage{physics}

\geometry{margin=1in}
\hypersetup{colorlinks=true, linkcolor=black, urlcolor=blue, citecolor=black}
\setlist[enumerate]{topsep=0.4em,itemsep=0.35em}
\setlist[itemize]{topsep=0.3em,itemsep=0.25em}
\titleformat{\section}{\large\bfseries}{\thesection.}{0.5em}{}
\titleformat{\subsection}{\normalsize\bfseries}{\thesubsection.}{0.5em}{}
\setstretch{1.06}

\newcommand{\Aether}{\AE{}ther}
\newcommand{\Aflow}{\AE{}ther-flow}
\newcommand{\TheoryName}{\Aflow{} Interpretation of Relativity}
\newcommand{\TheoryProgram}{\Aether{} / \Aflow{} framework}

\newcommand{\CompletedMetric}{g^{\sharp}}
\newcommand{\MatterSpace}{\Psi}

\newcommand{\Car}{\mathrm{car}}
\newcommand{\Cur}{\mathrm{cur}}
\newcommand{\Hom}{\mathrm{hom}}

\newcommand{\ForSpace}[1]{\mathcal I_{#1}^{\mathrm{for}}}
\newcommand{\PrimShell}{\mathcal S_{\mathrm{prim}}}
\newcommand{\Reduction}[1]{\mathcal R_{#1}}
\newcommand{\CurrentCarrier}[1]{\mathfrak C_{#1}^{\mathrm{cur}}}
\newcommand{\EqCarrier}[1]{\mathfrak C_{#1}^{\mathrm{eq}}}
\newcommand{\MetricOut}{\mathscr G_{\mathrm{out}}}
\newcommand{\MatterOut}{\mathscr T_{\mathrm{out}}}

\newcommand{\SubTarget}[1]{E_{#1}^{\mathrm{sub}}}
\newcommand{\SubPair}[1]{\mathfrak s_{#1}^{\mathrm{sub}}}
\newcommand{\EqnTarget}[1]{Q_{#1}^{\mathrm{eqn}}}
\newcommand{\HiddenSpace}[1]{\mathcal H_{#1}^{\mathrm{hid}}}
\newcommand{\HiddenBody}[1]{D_{#1}^{\mathrm{hid}}}

\newcommand{\ProtoSpace}[1]{\mathcal U_{#1}^{\mathrm{proto}}}
\newcommand{\ProtoBody}[1]{\mathcal B_{#1}^{\mathrm{proto}}}
\newcommand{\ProtoSection}[1]{\Gamma_{#1}^{\mathrm{proto}}}
\newcommand{\ProtoShellSet}[1]{\mathcal Z_{#1}^{\mathrm{proto}}}
\newcommand{\ProtoZeroCore}[1]{\mathcal Z_{#1}^{0,\mathrm{proto}}}
\newcommand{\ProtoSplit}[1]{\Phi_{#1}^{\mathrm{proto}}}
\newcommand{\ProtoCharge}[1]{\mathcal Q_{#1}^{\mathrm{proto}}}
\newcommand{\ProtoEmbed}[1]{\zeta_{#1}^{\mathrm{proto}}}
\newcommand{\ProtoKernelCh}[1]{\widetilde K_{#1}^{0,\mathrm{proto}}}
\newcommand{\ProtoStateComp}[1]{P_{#1}^{\mathrm{proto,co}}}

\newcommand{\PreProtoNucleus}{\mathfrak Q^{\mathrm{nuc}}}
\newcommand{\PreProtoSpace}[1]{\mathcal V_{#1}^{\mathrm{preproto}}}
\newcommand{\PreProtoBody}[1]{\mathcal B_{#1}^{\mathrm{preproto}}}
\newcommand{\PreProtoProj}[1]{\Pi_{#1}^{\mathrm{preproto}}}
\newcommand{\PreProtoProtoMin}[1]{\mathfrak f_{#1}^{\mathrm{preproto,proto}}}
\newcommand{\PreProtoSection}[1]{\Gamma_{#1}^{\mathrm{preproto}}}
\newcommand{\PreProtoFunctional}[1]{\mathscr W_{#1}^{\mathrm{preproto}}}
\newcommand{\PreProtoShellSet}[1]{\mathcal Z_{#1}^{\mathrm{preproto}}}

\newcommand{\PreNucCore}{\mathfrak R^{\mathrm{core}}}
\newcommand{\PreNucSpace}[1]{\mathcal W_{#1}^{\mathrm{prenuc}}}
\newcommand{\PreNucBody}[1]{\mathcal B_{#1}^{\mathrm{prenuc}}}
\newcommand{\PreNucProj}[1]{\Pi_{#1}^{\mathrm{prenuc}}}
\newcommand{\PreNucNucleusMin}[1]{\mathfrak f_{#1}^{\mathrm{prenuc,nuc}}}
\newcommand{\PreNucShellSet}[1]{\mathcal Z_{#1}^{\mathrm{prenuc}}}

\newcommand{\PreCorePackage}{\mathfrak S^{\mathrm{precore}}}
\newcommand{\PreCoreTransportCore}{\mathfrak S^{\mathrm{tr}}}
\newcommand{\PreCoreSpace}[1]{\mathcal Y_{#1}^{\mathrm{precore}}}
\newcommand{\PreCoreBody}[1]{\mathcal B_{#1}^{\mathrm{precore}}}
\newcommand{\PreCoreProj}[1]{\Pi_{#1}^{\mathrm{precore}}}
\newcommand{\PreCoreCoreMin}[1]{\mathfrak f_{#1}^{\mathrm{precore,core}}}
\newcommand{\PreCoreShellSet}[1]{\mathcal Z_{#1}^{\mathrm{precore}}}
\newcommand{\PreCoreTrBody}[1]{\mathcal B_{#1}^{\mathrm{precore,tr}}}
\newcommand{\PreCoreTrProj}[1]{\widehat{\Pi}_{#1}^{\mathrm{precore}}}
\newcommand{\PreCoreTrShell}[1]{\mathcal Z_{#1}^{\mathrm{precore,tr}}}

\newtheorem{definition}{Definition}
\newtheorem{proposition}{Proposition}
\newtheorem{remark}{Remark}
\newtheorem{corollary}{Corollary}
\newtheorem{theorem}{Theorem}

\title{\textbf{The \TheoryName{}}\\[0.4em]
\large Primitive-Reservoir Observer-Reduction\\
\large Exact-Shell-Core-Generating Substrate Pre-Core Dynamics\\
\large Collapse to an Observer-Relevant Pre-Core Exact-Shell Transport Core}
\author{}
\date{}

\begin{document}

\maketitle

\begin{abstract}
The current primitive-reservoir observer-reduction frontier already removed the observer-relevant pre-nucleus exact-shell core as a primitive stopping point by deriving it from exact-shell-core-generating substrate pre-core dynamics. After that step, the honest next move is no longer another same-output deeper-origin relay unless it changes the benchmark-facing burden. At this stage there are only two admissible routes: derive or justify the pre-core dynamics themselves from still more primitive substrate structure, or prove that the full pre-core package still contains benchmark-neutral ambient freedom and collapses to a smaller theorem-level datum.

The present note takes the second route. It proves that the full pre-core package carries no independent benchmark-facing freedom beyond its minimizer image over the recorded pre-nucleus exact-shell core. More precisely, what remains observer relevant at this frontier is the \emph{pre-core exact-shell transport core}: the minimizer-image compact transport body over the recorded pre-nucleus body, the restricted projection to that body, the exact-shell transport image of the recorded pre-nucleus shell, and benchmark faithfulness. The full exact pre-core shell is forced to coincide with that transport image, and ambient pre-core directions outside the minimizer image do not add new benchmark-facing data once the downstream exact-shell-core target is held fixed.

The benchmark boundary remains fixed throughout. The one operative metric is still exactly \(\CompletedMetric\), the ordinary matter route is unchanged, there is no second operative metric, the low-energy matter law is not enlarged, and no stronger promotion language is licensed. The positive result is therefore narrow but benchmark facing: the live burden below the current pre-core frontier collapses from the full exact-shell-core-generating substrate pre-core package \(\PreCorePackage\) to the smaller observer-relevant pre-core exact-shell transport core \(\PreCoreTransportCore\). What remains open is to derive or justify that sharper transport core from still more primitive substrate structure, or to prove an even smaller collapse strictly inside it.
\end{abstract}

\section{Introduction}

The active derivational program of \TheoryName{} is now disciplined enough that each new manuscript must either move the live burden downward or compress it further. The current active-primary theorem already took the first route at pre-core scope: it showed that the observer-relevant pre-nucleus exact-shell core is the projected exact-shell transport image of a deeper pre-core package. That result removed the former exact-shell-core stopping point, but it did not yet answer whether the full pre-core package is itself already minimal.

The present note takes the remaining internal-compression route at that frontier. The question is no longer whether the current pre-core package has some presentation at all; it does. The sharper question is whether every part of that package is still benchmark facing once the recorded pre-nucleus exact-shell core and the benchmark one-metric package are held fixed. The theorem proved here answers that question conservatively. The full pre-core package still contains ambient directions that do not add independent benchmark-facing freedom.

What survives as the live burden is smaller. Only the minimizer-image transport body over the recorded pre-nucleus body, the restricted projection back to that body, the exact-shell transport image, and benchmark faithfulness matter at benchmark scope. The ambient pre-core state space and off-image body directions remain part of one possible presentation, but they do not survive as additional observer-relevant freedom once the exact-shell transport image is fixed.

\paragraph{Framework and claim boundary.}
\TheoryProgram{} denotes the broader \Aether{} / \Aflow{} framework built on the claims that the \Aether{} is the underlying four-dimensional substrate of reality, the \Aflow{} is its intrinsic ordered motion, observed three-dimensional space is the local experiential slice of that deeper substrate, S-time is the experienced order of change arising from matter, light, and the \Aflow{}, and observed expansion is the three-dimensional appearance of deeper four-dimensional ordered motion. Within that framework, \TheoryName{} denotes the exact-closure theory in which the effective gravitational dynamics are adopted to be exactly Einsteinian with universal matter coupling. In the active sequence, \emph{adoption} means use of that established relativistic dynamics without claiming substrate derivation, while \emph{derivation} is reserved for a first-principles recovery from explicit substrate variables.

The present note stays strictly inside that boundary. It does not alter the benchmark exact-closure package. It does not introduce a second operative metric or a new low-energy matter law. It only sharpens what now counts as the benchmark-facing burden inside the current pre-core frontier.

\section{Fixed Inputs}

\begin{definition}[Recorded primitive continuation data]
\label{def:fixed}
Fix the following continuation data from the active primitive-reservoir observer-reduction line.
\begin{enumerate}
    \item For each sector label \(\sigma\in\{\Car,\Cur,\Hom\}\), a primitive forcing shell
    \[
    \ForSpace{\sigma},
    \]
    a visible reduction map
    \[
    \Reduction{\sigma}:\ForSpace{\sigma}\to X_{\sigma},
    \]
    and shell-level current and equilibrium carriers
    \[
    \CurrentCarrier{\sigma}:\ForSpace{\sigma}\to Z_{\sigma}^{\mathrm{cur}},
    \qquad
    \EqCarrier{\sigma}:\ForSpace{\sigma}\to Z_{\sigma}^{\mathrm{eq}}.
    \]
    \item The product shell
    \[
    \PrimShell
    \equiv
    \ForSpace{\Car}\times \ForSpace{\Cur}\times \ForSpace{\Hom}.
    \]
    \item The recorded metric and ordinary-matter outputs
    \[
    \MetricOut:\PrimShell\to\{\text{observer metrics}\},
    \qquad
    \MatterOut:\PrimShell\times\MatterSpace\to\{\text{observer stress tensors}\},
    \]
    satisfying
    \begin{equation}
    \MetricOut(\mathbf w)=\CompletedMetric,
    \qquad
    \MatterOut(\mathbf w,\psi)
    =
    T_{\mu\nu}^{\mathrm{matter}}[\CompletedMetric,\psi]
    \label{eq:fixedoutputs}
    \end{equation}
    for every \(\mathbf w\in\PrimShell\) and every matter configuration \(\psi\in\MatterSpace\).
\end{enumerate}
\end{definition}

\begin{definition}[Recorded proto-germ exact-shell target]
\label{def:proto}
Fix \(\sigma\in\{\Car,\Cur,\Hom\}\). The active line already fixes:
\begin{enumerate}
    \item the same substrate target and pairing
    \[
    \SubTarget{\sigma},
    \qquad
    \SubPair{\sigma}:
    \SubTarget{\sigma}\times\SubTarget{\sigma}\to\mathbb R;
    \]
    \item a hidden fiber space \(\HiddenSpace{\sigma}\) with compact hidden body \(\HiddenBody{\sigma}\subset\HiddenSpace{\sigma}\);
    \item a finite-dimensional proto-germ state space \(\ProtoSpace{\sigma}\), a compact convex proto-germ body \(\ProtoBody{\sigma}\subset\ProtoSpace{\sigma}\), a smooth proto-germ restriction section
    \[
    \ProtoSection{\sigma}:\ProtoSpace{\sigma}\to\SubTarget{\sigma},
    \]
    and the exact proto-germ shell
    \[
    \ProtoShellSet{\sigma}
    \equiv
    \bigl(\ProtoSection{\sigma}\bigr)^{-1}(0)\cap\ProtoBody{\sigma};
    \]
    \item a closed embedded proto zero core \(\ProtoZeroCore{\sigma}\subset\ProtoShellSet{\sigma}\) and a split diffeomorphism
    \[
    \ProtoSplit{\sigma}:
    \ProtoZeroCore{\sigma}\times\HiddenBody{\sigma}
    \xrightarrow{\cong}
    \ProtoShellSet{\sigma};
    \]
    \item a hidden charge
    \[
    \ProtoCharge{\sigma}:\HiddenSpace{\sigma}\to\EqnTarget{\sigma},
    \]
    a smooth observer-compatible exact proto-germ zero-response realization
    \[
    \ProtoEmbed{\sigma}:\ForSpace{\sigma}\hookrightarrow\ProtoZeroCore{\sigma},
    \]
    and a fixed-data solvability / coercive package on \(\ProtoZeroCore{\sigma}\), recorded through \(\ProtoKernelCh{\sigma}\) and \(\ProtoStateComp{\sigma}\).
\end{enumerate}
\end{definition}

\begin{definition}[Recorded current transport nucleus]
\label{def:nucleus}
Fix \(\sigma\in\{\Car,\Cur,\Hom\}\). The recorded current transport nucleus \(\PreProtoNucleus\) for sector \(\sigma\) consists of:
\begin{enumerate}
    \item the pre-proto-germ state space \(\PreProtoSpace{\sigma}\), compact convex body \(\PreProtoBody{\sigma}\subset\PreProtoSpace{\sigma}\), projection
    \[
    \PreProtoProj{\sigma}:\PreProtoSpace{\sigma}\to\ProtoSpace{\sigma},
    \]
    and proto section
    \[
    \PreProtoProtoMin{\sigma}:\ProtoBody{\sigma}\to\PreProtoBody{\sigma}
    \]
    satisfying
    \[
    \PreProtoProj{\sigma}\bigl(\PreProtoBody{\sigma}\bigr)=\ProtoBody{\sigma},
    \qquad
    \PreProtoProj{\sigma}\circ\PreProtoProtoMin{\sigma}
    =
    \mathrm{id}_{\ProtoBody{\sigma}};
    \]
    \item the restriction section
    \[
    \PreProtoSection{\sigma}:\PreProtoSpace{\sigma}\to\SubTarget{\sigma},
    \]
    the induced functional
    \[
    \PreProtoFunctional{\sigma}(q)
    \equiv
    \SubPair{\sigma}\bigl(\PreProtoSection{\sigma}(q),\PreProtoSection{\sigma}(q)\bigr),
    \]
    and the exact shell
    \[
    \PreProtoShellSet{\sigma}
    \equiv
    \bigl(\PreProtoSection{\sigma}\bigr)^{-1}(0)\cap\PreProtoBody{\sigma},
    \]
    with
    \[
    \PreProtoSection{\sigma}\circ\PreProtoProtoMin{\sigma}
    =
    \ProtoSection{\sigma},
    \qquad
    \PreProtoProj{\sigma}\big|_{\PreProtoShellSet{\sigma}}
    :
    \PreProtoShellSet{\sigma}
    \xrightarrow{\cong}
    \ProtoShellSet{\sigma};
    \]
    \item benchmark faithfulness, meaning preservation of the benchmark outputs \eqref{eq:fixedoutputs}, no second operative metric, no enlarged low-energy matter law, and no premature promotion from adoption to completed derivation.
\end{enumerate}
\end{definition}

\begin{definition}[Recorded observer-relevant pre-nucleus exact-shell core]
\label{def:core}
Fix \(\sigma\in\{\Car,\Cur,\Hom\}\). The observer-relevant pre-nucleus exact-shell core \(\PreNucCore\) for sector \(\sigma\) consists of:
\begin{enumerate}
    \item the pre-nucleus state space \(\PreNucSpace{\sigma}\), compact convex body \(\PreNucBody{\sigma}\subset\PreNucSpace{\sigma}\), affine surjection
    \[
    \PreNucProj{\sigma}:\PreNucSpace{\sigma}\to\PreProtoSpace{\sigma},
    \]
    and fiber minimizer
    \[
    \PreNucNucleusMin{\sigma}:\PreProtoBody{\sigma}\to\PreNucBody{\sigma}
    \]
    satisfying
    \[
    \PreNucProj{\sigma}\bigl(\PreNucBody{\sigma}\bigr)=\PreProtoBody{\sigma},
    \qquad
    \PreNucProj{\sigma}\circ\PreNucNucleusMin{\sigma}
    =
    \mathrm{id}_{\PreProtoBody{\sigma}};
    \]
    \item a closed embedded exact pre-nucleus shell
    \[
    \PreNucShellSet{\sigma}\subset\PreNucBody{\sigma}
    \]
    satisfying
    \[
    \PreNucNucleusMin{\sigma}\bigl(\PreProtoShellSet{\sigma}\bigr)
    \subset
    \PreNucShellSet{\sigma},
    \qquad
    \PreNucProj{\sigma}\big|_{\PreNucShellSet{\sigma}}
    :
    \PreNucShellSet{\sigma}
    \xrightarrow{\cong}
    \PreProtoShellSet{\sigma};
    \]
    \item benchmark faithfulness in the sense of Definition \ref{def:nucleus}.
\end{enumerate}
\end{definition}

\begin{definition}[Exact-shell-core-generating substrate pre-core dynamics]
\label{def:precore}
Fix \(\sigma\in\{\Car,\Cur,\Hom\}\). An exact-shell-core-generating substrate pre-core dynamics package \(\PreCorePackage\) for sector \(\sigma\) consists of the following data.
\begin{enumerate}
    \item A finite-dimensional Euclidean pre-core state space \(\PreCoreSpace{\sigma}\), a compact convex pre-core body
    \[
    \PreCoreBody{\sigma}\subset\PreCoreSpace{\sigma}
    \]
    with nonempty interior, an affine surjection
    \[
    \PreCoreProj{\sigma}:\PreCoreSpace{\sigma}\to\PreNucSpace{\sigma},
    \]
    and a smooth pre-core minimizer
    \[
    \PreCoreCoreMin{\sigma}:\PreNucBody{\sigma}\to\PreCoreBody{\sigma}
    \]
    satisfying
    \[
    \PreCoreProj{\sigma}\bigl(\PreCoreBody{\sigma}\bigr)=\PreNucBody{\sigma},
    \qquad
    \PreCoreProj{\sigma}\circ\PreCoreCoreMin{\sigma}
    =
    \mathrm{id}_{\PreNucBody{\sigma}}.
    \]
    \item A closed embedded exact pre-core shell
    \[
    \PreCoreShellSet{\sigma}\subset\PreCoreBody{\sigma}
    \]
    satisfying
    \[
    \PreCoreCoreMin{\sigma}\bigl(\PreNucShellSet{\sigma}\bigr)
    \subset
    \PreCoreShellSet{\sigma},
    \qquad
    \PreCoreProj{\sigma}\big|_{\PreCoreShellSet{\sigma}}
    :
    \PreCoreShellSet{\sigma}
    \xrightarrow{\cong}
    \PreNucShellSet{\sigma}.
    \]
    \item Benchmark faithfulness, meaning preservation of the benchmark outputs \eqref{eq:fixedoutputs}, no second operative metric, no enlarged low-energy matter law, and no premature promotion from adoption to completed derivation.
\end{enumerate}
\end{definition}

\begin{remark}
Definitions \ref{def:fixed}--\ref{def:precore} are fixed inputs. The one operative metric, the ordinary matter route, the fixed proto-germ exact-shell target, the current transport nucleus, the recorded pre-nucleus exact-shell core, and the existence of a benchmark-faithful pre-core package are not reopened here. The present note asks only whether the full pre-core package is already minimal once those downstream targets are held fixed.
\end{remark}

\section{Observer-Relevant Pre-Core Exact-Shell Transport Core}

\begin{definition}[Observer-relevant pre-core exact-shell transport core]
\label{def:transportcore}
Fix \(\sigma\in\{\Car,\Cur,\Hom\}\) and a benchmark-faithful pre-core package in the sense of Definition \ref{def:precore}. The \emph{observer-relevant pre-core exact-shell transport core} \(\PreCoreTransportCore\) for sector \(\sigma\) consists of:
\begin{enumerate}
    \item the compact transport-image body
    \[
    \PreCoreTrBody{\sigma}
    \equiv
    \PreCoreCoreMin{\sigma}\bigl(\PreNucBody{\sigma}\bigr)
    \subset
    \PreCoreBody{\sigma},
    \]
    together with the restricted projection
    \[
    \PreCoreTrProj{\sigma}
    \equiv
    \PreCoreProj{\sigma}\big|_{\PreCoreTrBody{\sigma}}
    :
    \PreCoreTrBody{\sigma}\to\PreNucBody{\sigma};
    \]
    \item the exact-shell transport image
    \[
    \PreCoreTrShell{\sigma}
    \equiv
    \PreCoreCoreMin{\sigma}\bigl(\PreNucShellSet{\sigma}\bigr)
    \subset
    \PreCoreShellSet{\sigma},
    \]
    with the restricted shell projection
    \[
    \PreCoreTrProj{\sigma}\big|_{\PreCoreTrShell{\sigma}}
    :
    \PreCoreTrShell{\sigma}\to\PreNucShellSet{\sigma};
    \]
    \item benchmark faithfulness in the sense of Definition \ref{def:precore}.
\end{enumerate}
\end{definition}

\begin{remark}
Definition \ref{def:transportcore} is intentionally smaller than Definition \ref{def:precore}. It keeps only the minimizer image over the recorded pre-nucleus exact-shell core together with the exact-shell transport it already determines. The ambient pre-core state space and off-image body directions are not taken as independent core data.
\end{remark}

\section{Collapse of the Full Pre-Core Package}

\begin{proposition}[The transport-image body carries the body-level content]
\label{prop:body}
For every benchmark-faithful pre-core package, \(\PreCoreTrBody{\sigma}\) is compact and the restricted projection
\[
\PreCoreTrProj{\sigma}:\PreCoreTrBody{\sigma}\to\PreNucBody{\sigma}
\]
is a homeomorphism with inverse \(\PreCoreCoreMin{\sigma}\).
\end{proposition}

\begin{proof}
The body \(\PreNucBody{\sigma}\) is compact by Definition \ref{def:core}, and \(\PreCoreCoreMin{\sigma}\) is continuous by Definition \ref{def:precore}(1). Therefore
\[
\PreCoreTrBody{\sigma}
=
\PreCoreCoreMin{\sigma}\bigl(\PreNucBody{\sigma}\bigr)
\]
is compact.

By Definition \ref{def:precore}(1),
\[
\PreCoreProj{\sigma}\circ\PreCoreCoreMin{\sigma}
=
\mathrm{id}_{\PreNucBody{\sigma}}.
\]
Hence
\[
\PreCoreTrProj{\sigma}\circ\PreCoreCoreMin{\sigma}
=
\mathrm{id}_{\PreNucBody{\sigma}}.
\]
Conversely, every \(y\in\PreCoreTrBody{\sigma}\) has the form
\[
y=\PreCoreCoreMin{\sigma}(w)
\qquad
\text{for some }
w\in\PreNucBody{\sigma},
\]
so
\begin{align*}
\PreCoreCoreMin{\sigma}\bigl(\PreCoreTrProj{\sigma}(y)\bigr)
&=
\PreCoreCoreMin{\sigma}\bigl(\PreCoreProj{\sigma}(y)\bigr) \\
&=
\PreCoreCoreMin{\sigma}\bigl(\PreCoreProj{\sigma}(\PreCoreCoreMin{\sigma}(w))\bigr) \\
&=
\PreCoreCoreMin{\sigma}(w) \\
&=
y.
\end{align*}
Thus \(\PreCoreTrProj{\sigma}\) and \(\PreCoreCoreMin{\sigma}\) are mutually inverse continuous maps between compact Hausdorff spaces, hence homeomorphisms.
\end{proof}

\begin{proposition}[The full exact pre-core shell is forced by the transport image]
\label{prop:shell}
For every benchmark-faithful pre-core package,
\[
\PreCoreTrShell{\sigma}
=
\PreCoreShellSet{\sigma},
\]
and the restricted shell projection
\[
\PreCoreTrProj{\sigma}\big|_{\PreCoreTrShell{\sigma}}
:
\PreCoreTrShell{\sigma}\xrightarrow{\cong}\PreNucShellSet{\sigma}
\]
is a diffeomorphism with inverse \(\PreCoreCoreMin{\sigma}\big|_{\PreNucShellSet{\sigma}}\).
\end{proposition}

\begin{proof}
Definition \ref{def:precore}(2) already gives the inclusion
\[
\PreCoreCoreMin{\sigma}\bigl(\PreNucShellSet{\sigma}\bigr)
\subset
\PreCoreShellSet{\sigma},
\]
which is the inclusion
\[
\PreCoreTrShell{\sigma}\subset\PreCoreShellSet{\sigma}.
\]

For the reverse inclusion, let \(z\in\PreCoreShellSet{\sigma}\). Since
\[
\PreCoreProj{\sigma}\big|_{\PreCoreShellSet{\sigma}}
:
\PreCoreShellSet{\sigma}\xrightarrow{\cong}\PreNucShellSet{\sigma}
\]
is a diffeomorphism by Definition \ref{def:precore}(2), the point
\[
w\equiv\PreCoreProj{\sigma}(z)
\]
lies in \(\PreNucShellSet{\sigma}\). Again by Definition \ref{def:precore}(2),
\[
\PreCoreCoreMin{\sigma}(w)\in\PreCoreShellSet{\sigma}.
\]
But
\[
\PreCoreProj{\sigma}\bigl(\PreCoreCoreMin{\sigma}(w)\bigr)=w=\PreCoreProj{\sigma}(z),
\]
and the shell projection is injective because it is a diffeomorphism. Therefore
\[
z=\PreCoreCoreMin{\sigma}(w)\in\PreCoreTrShell{\sigma}.
\]
Hence
\[
\PreCoreShellSet{\sigma}\subset\PreCoreTrShell{\sigma},
\]
so equality holds. The final statement follows because the full shell projection is already a diffeomorphism and its inverse on \(\PreNucShellSet{\sigma}\) is exactly \(\PreCoreCoreMin{\sigma}\big|_{\PreNucShellSet{\sigma}}\).
\end{proof}

\begin{proposition}[Downstream exact-shell transport factors through the transport core]
\label{prop:downstream}
For every benchmark-faithful pre-core package, the recovery of the recorded pre-nucleus exact-shell core and all downstream shell transport to the fixed transport nucleus and proto-germ exact-shell target factor through \(\PreCoreTransportCore\). More precisely,
\[
\PreNucProj{\sigma}\circ\PreCoreTrProj{\sigma}
:
\PreCoreTrBody{\sigma}\to\PreProtoBody{\sigma}
\]
is a surjection with section
\[
\PreCoreCoreMin{\sigma}\circ\PreNucNucleusMin{\sigma}
:
\PreProtoBody{\sigma}\to\PreCoreTrBody{\sigma},
\]
and the composite shell transport
\[
(\PreNucProj{\sigma}\circ\PreCoreTrProj{\sigma})\big|_{\PreCoreTrShell{\sigma}}
:
\PreCoreTrShell{\sigma}
\xrightarrow{\cong}
\PreProtoShellSet{\sigma}
\]
is a diffeomorphism. Consequently
\[
(\PreProtoProj{\sigma}\circ\PreNucProj{\sigma}\circ\PreCoreTrProj{\sigma})\big|_{\PreCoreTrShell{\sigma}}
:
\PreCoreTrShell{\sigma}
\xrightarrow{\cong}
\ProtoShellSet{\sigma}
\]
is also a diffeomorphism.
\end{proposition}

\begin{proof}
By Proposition \ref{prop:body},
\[
\PreCoreTrProj{\sigma}\circ\PreCoreCoreMin{\sigma}
=
\mathrm{id}_{\PreNucBody{\sigma}}.
\]
Composing with Definition \ref{def:core}(1) yields
\[
(\PreNucProj{\sigma}\circ\PreCoreTrProj{\sigma})\bigl(\PreCoreTrBody{\sigma}\bigr)
=
\PreNucProj{\sigma}\bigl(\PreNucBody{\sigma}\bigr)
=
\PreProtoBody{\sigma}
\]
and
\[
(\PreNucProj{\sigma}\circ\PreCoreTrProj{\sigma})\circ
\PreCoreCoreMin{\sigma}\circ\PreNucNucleusMin{\sigma}
=
\PreNucProj{\sigma}\circ\PreNucNucleusMin{\sigma}
=
\mathrm{id}_{\PreProtoBody{\sigma}}.
\]

For shell transport, Proposition \ref{prop:shell} gives a diffeomorphism
\[
\PreCoreTrProj{\sigma}\big|_{\PreCoreTrShell{\sigma}}
:
\PreCoreTrShell{\sigma}\xrightarrow{\cong}\PreNucShellSet{\sigma},
\]
and Definition \ref{def:core}(2) gives a diffeomorphism
\[
\PreNucProj{\sigma}\big|_{\PreNucShellSet{\sigma}}
:
\PreNucShellSet{\sigma}\xrightarrow{\cong}\PreProtoShellSet{\sigma}.
\]
Their composition is therefore a diffeomorphism from \(\PreCoreTrShell{\sigma}\) to \(\PreProtoShellSet{\sigma}\). Composing once more with the shell diffeomorphism of Definition \ref{def:nucleus}(2) yields the final displayed diffeomorphism to \(\ProtoShellSet{\sigma}\).
\end{proof}

\begin{proposition}[Ambient pre-core directions are benchmark neutral at current scope]
\label{prop:neutral}
For every benchmark-faithful pre-core package, the ambient pre-core state space \(\PreCoreSpace{\sigma}\) and the off-image body directions
\[
\PreCoreBody{\sigma}\setminus\PreCoreTrBody{\sigma}
\]
do not add independent benchmark-facing freedom beyond \(\PreCoreTransportCore\).
\end{proposition}

\begin{proof}
By Proposition \ref{prop:body}, the full body-level transport relevant to the recorded pre-nucleus core is already fixed by the transport-image body \(\PreCoreTrBody{\sigma}\), the restricted projection \(\PreCoreTrProj{\sigma}\), and the minimizer \(\PreCoreCoreMin{\sigma}\). By Proposition \ref{prop:shell}, the full exact pre-core shell is not additional data beyond the transport image of the recorded pre-nucleus shell; it is exactly that image. By Proposition \ref{prop:downstream}, all downstream transport to the recorded transport nucleus and the fixed proto-germ exact-shell target factors through this smaller transport core. Finally, benchmark faithfulness is already part of Definition \ref{def:transportcore}(3) and preserves the fixed outputs \eqref{eq:fixedoutputs}. Therefore nothing benchmark-facing at the present frontier requires the ambient pre-core state-space presentation or body points outside the minimizer image as additional independent data.
\end{proof}

\begin{proposition}[The benchmark package remains fixed]
\label{prop:benchmark}
Every observer-relevant pre-core exact-shell transport core preserves the same benchmark one-metric package \eqref{eq:fixedoutputs}. In particular, the operative metric remains exactly \(\CompletedMetric\), the ordinary matter route is unchanged, there is no second operative metric, and the low-energy matter law is not enlarged.
\end{proposition}

\begin{proof}
This is part of benchmark faithfulness in Definition \ref{def:transportcore}(3). The present note only removes benchmark-neutral ambient pre-core freedom from the live burden. It does not alter the benchmark exact-closure package.
\end{proof}

\section{Main Theorem}

\begin{theorem}[Exact-shell-core-generating substrate pre-core dynamics collapse to an observer-relevant pre-core exact-shell transport core]
\label{thm:main}
Fix the primitive continuation data of Definition \ref{def:fixed}, the recorded proto-germ exact-shell target of Definition \ref{def:proto}, the recorded current transport nucleus of Definition \ref{def:nucleus}, the recorded observer-relevant pre-nucleus exact-shell core of Definition \ref{def:core}, and a benchmark-faithful pre-core package in the sense of Definition \ref{def:precore}. Then, for each sector \(\sigma\in\{\Car,\Cur,\Hom\}\), the live benchmark-facing burden inside the current pre-core frontier collapses from the full pre-core package \(\PreCorePackage\) to the smaller observer-relevant pre-core exact-shell transport core \(\PreCoreTransportCore\). More precisely:
\begin{enumerate}
    \item the body-level content of the full pre-core package is already carried by the transport-image body \(\PreCoreTrBody{\sigma}\) and the restricted projection \(\PreCoreTrProj{\sigma}\) by Proposition \ref{prop:body};
    \item the full exact pre-core shell is forced to equal the transport image of the recorded pre-nucleus shell by Proposition \ref{prop:shell};
    \item all downstream recovery of the recorded transport nucleus and fixed proto-germ exact-shell target factors through \(\PreCoreTransportCore\) by Proposition \ref{prop:downstream};
    \item ambient pre-core directions outside that transport core are benchmark neutral at current scope by Proposition \ref{prop:neutral};
    \item the benchmark boundary remains fixed by Proposition \ref{prop:benchmark};
    \item consequently the remaining honest burden below the previous pre-core frontier is no longer the full exact-shell-core-generating substrate pre-core dynamics package \(\PreCorePackage\) as such. What remains is to derive or justify the sharper transport core \(\PreCoreTransportCore\) from still more primitive substrate structure, or else prove a still smaller theorem-level collapse strictly inside that transport core.
\end{enumerate}
\end{theorem}

\begin{proof}
Item (1) is Proposition \ref{prop:body}. Item (2) is Proposition \ref{prop:shell}. Item (3) is Proposition \ref{prop:downstream}. Item (4) is Proposition \ref{prop:neutral}. Item (5) is Proposition \ref{prop:benchmark}. Item (6) is the routing consequence of the previous items together with the active safeguard against counting another same-output deeper-origin relay as new front-line progress when the live benchmark-relevant datum is unchanged.
\end{proof}

\begin{corollary}[What must be done next]
\label{cor:next}
After Theorem \ref{thm:main}, the next honest benchmark-facing manuscript must either:
\begin{enumerate}
    \item derive or justify the observer-relevant pre-core exact-shell transport core \(\PreCoreTransportCore\) from still more primitive substrate structure in a way that changes benchmark-facing status;
    \item identify a genuinely smaller theorem-level observer-relevant datum strictly inside \(\PreCoreTransportCore\) and prove a sharper collapse to it;
    \item do either of the previous two while preserving the same benchmark one-metric package and the same claim boundary.
\end{enumerate}
Another same-output deeper-origin relay that preserves this transport core, another re-expansion back to the full pre-core package as if its screened ambient directions were still live, another reopening of the former pre-nucleus exact-shell-core frontier as if it were again the active burden, or any move that introduces a second operative metric, enlarges the ordinary matter law, or uses stronger promotion language does not satisfy the present burden. If a quantum continuation is pursued, it matters benchmark-facingly only insofar as it derives this same transport core, collapses it further, or closes a benchmark derivation gate in a genuinely new way.
\end{corollary}

\begin{proof}
Theorem \ref{thm:main} shows that the full pre-core package is no longer the minimal benchmark-facing datum at this stage. Therefore a subsequent manuscript must improve on the smaller transport core rather than merely restating the larger package.
\end{proof}

\section{Conclusion}

The positive result of this note is a disciplined compression step below the current pre-core frontier. The previous active-primary theorem already removed the observer-relevant pre-nucleus exact-shell core as the primitive stopping point by deriving it from a deeper pre-core package. The present theorem then takes the remaining honest internal route at that scope: it shows that the full pre-core package is itself not the minimal benchmark-facing datum.

What survives as the live burden is the observer-relevant pre-core exact-shell transport core. That sharper datum consists of the minimizer-image compact transport body over the recorded pre-nucleus body, the restricted projection back to that body, the exact-shell transport image of the recorded pre-nucleus shell, and benchmark faithfulness. The full exact pre-core shell is forced to equal that transport image, and ambient pre-core directions outside the minimizer image do not add independent benchmark-facing freedom once the downstream exact-shell-core target is fixed.

This is the honest next burden after the present theorem. A new primary manuscript must derive or justify that sharper transport core from still more primitive substrate structure, or else collapse it further, while preserving the same benchmark one-metric package and the same claim boundary. The current result does not yet complete a first-principles derivation of GR from substrate microphysics, but it does narrow the live derivational burden one level further on the active line.

\end{document}
